%%
%% fall-lecture14.tex
%% 
%% Made by Alex Nelson <pqnelson@gmail.com>
%% Login   <alex@lisp>
%% 
%% Started on  2026-11-03T10:03:40-0800
%% Last update 2026-11-03T10:03:40-0800
%% 

\lecture{}

\begin{definition}
Let $H$ be a Hilbert space. Let $y\in H$ be a vector.
Then $f\colon H\to\CC$, $x\mapsto(x,y)$ is a continuous linear map $f\in\mathcal{L}(H,\CC)$.
We define the \define{Dual Space} of $H$ to be $H' := \mathcal{L}(H,\CC)$.
\end{definition}

\begin{theorem}[Riesz representation]
Let $H$ be a Hilbert space. 
For every element of the dual space $f\in H'$,
there exists a unique $y\in H$ such that for all $x\in H$ we have $f(x)=(x,y)$.
\end{theorem}

\begin{proof}
Write $L=\ker(f)$, which is a closed linear subspace of $H$ (since $f$
is continuous). We have two cases:

\textsc{Case 1}: $L=H$. Then take $y=0$, so $f(x)=(x,y)$ for all $x\in H$.

\textsc{Case 2}: $L\neq H$. Then take any $y_{0}\in L^{\perp}$ of norm
$1=\|y_{0}\|$. Recall the Riesz projection Theorem~\ref{thm:riesz-projection}
which tells us $H=L\oplus L^{\perp}$. We will show that
\begin{equation}
y = \overline{f(y_{0})}y_{0}
\end{equation}
is the unique vector we seek. Observe $y\in L^{\perp}$.

We will show that for all $x\in H$, we have $f(x)=(x,y)$. We have two
subcases:

\textsc{Subcase 2.1}: $x\in L$. Then $f(x)=0=(x,y)$ since $x\in L$ but
$y\in L^{\perp}$.

\textsc{Subcase 2.2}: $x=y_{0}$. Then
\begin{subequations}
  \begin{align}
(x,y) &= f(y_{0})\cdot(y_{0},y_{0})\quad\mbox{by antilinearity in second slot}\\
&= f(y_{0})\|y_{0}\|^{2}\\
&= f(y_{0})\cdot1=f(y_{0})\\
&=f(x)\quad\mbox{by assumption}.
  \end{align}
\end{subequations}
Hence $f(x)=f(y_{0})=(x,y)$.

We claim $L^{\perp}=\CC y_{0}=\Span(y_{0})$. Then for $x\in L^{\perp}$
we can write it as $x=\lambda y_{0}$, and so:
\begin{equation}
f(x)=f(\lambda y_{0})=\lambda f(y_{0})=\lambda,
\end{equation}
which handles this case.

Hence $f(x)=(x,y)$ for all $x\in L\oplus L^{\perp}$ by these two
subcases, hence for all $x\in H$.

We just need to prove the claim $L^{\perp}=\CC y_{0}$: Let $x\in L^{\perp}$.
We know $\CC y_{0}\subset L^{\perp}$. We just need to prove $\CC y_{0}\supset L^{\perp}$.
Find nontrivial $\alpha,\beta\in\CC\setminus\{0\}$ such that
\begin{equation}
\alpha f(x)+\beta f(y_{0})=0.
\end{equation}
We can write
\begin{equation}
f(\alpha x + \beta y_{0})=\alpha f(x)+\beta f(y_{0})=0,
\end{equation}
so this means $\alpha x + \beta y_{0}\in L$. But both $x\in L^{\perp}$
and $y_{0}\in L^{\perp}$ implies $\alpha x + \beta y_{0}\in L^{\perp}$.
But the only vector which belongs to both $L$ and $L^{\perp}$ is the
zero vector, which means
\begin{equation}
\alpha x + \beta y_{0} = 0.
\end{equation}
This means any nonzero $x\in L^{\perp}$ is a nonzero scalar multiple of $y_{0}$.
Hence the inclusion $\CC y_{0}\supset L^{\perp}$ which concludes the claim.
\end{proof}

\subsection{Orthogonal Systems, Orthonormal Families}

\begin{definition}
Let $H$ be a Hilbert space. Let $e=\{e_{\alpha}\}_{\alpha\in A}\subset H$
be a family of vectors [possibly finite or not, possible uncountable or not].
\begin{itemize}
\item We say $e$ is an \define{Orthonormal Family} if
  $(e_{\alpha},e_{\beta})=\delta_{\alpha,\beta}$ for all
  $\alpha,\beta\in A$;
\item We say an orthonormal family $e$ is an \define{Orthonormal Basis}
  (or ``ONB'' for short) if one of the following equivalent conditions holds: 
\begin{enumerate}
\item all finite linear combinations $\Lin\{e_{\alpha}\}$ is dense in $H$
\item if $\forall\alpha\in A\ldotp(e_{\alpha},y)=0$, then $y=0$ [there
  is no nonzero vector orthogonal to the family]
\item the family $\{e_{\alpha}\}_{\alpha\in A}$ is maximal [i.e., we
  cannot add more orthogonal vectors to it]---formally: for any
  orthonormal family $e'$ such that $e\subset e'$ we have $e=e'$.
\end{enumerate}
\end{itemize}
\end{definition}

\begin{definition}
Let $H$ be a Hilbert space, let $\{e_{\alpha}\}$ be an orthonormal
basis. Let $x\in H$ be a vector. Then the \define{Fourier Expansion}
of $x$ with respect to $\{e_{\alpha}\}$ is the series $\sum_{\alpha\in A}\widehat{x}(\alpha)e_{\alpha}$.
(The sum is understood in terms of convergence of partial sums
directed by finite subsets; this is a directed set [net] directed by
finite subsets.)
We call $\widehat{x}(\alpha)$ the \define{Fourier Coefficients} for
$x$ with respect to the basis $\{e_{\alpha}\}$. 

It is true that $x=\sum_{\alpha\in A}\widehat{x}(\alpha)e_{\alpha}$.
\end{definition}

\begin{node}
Fourier expansions are studied further in 110C.
\end{node}

\begin{node}
The really difficult question concerns the \emph{existence} of an
orthonormal basis.
\end{node}

\begin{theorem}
Every Hilbert space has an orthonormal basis.
\end{theorem}

\begin{proof}
Zorn's lemma.
\end{proof}

\subsection{Aside: Zorn's Lemma}

\begin{definition}
A \define{Partially Ordered Set} (or ``poset'') is a set $A$ equipped
with a binary relation $\leq$ such that
\begin{enumerate}
\item Transitivity: for all $x,y,z\in A$, if $x\leq y$ and $y\leq z$, then $x\leq z$;
\item Reflexivity: for all $x\in A$, we have $x\leq x$.
\end{enumerate}
\end{definition}

\begin{definition}
Let $(A,\leq)$ be a poset. Let $E\subset A$ be a subset. We say $E$ is
\define{Linearly Ordered} if for all $a,b\in E$ either $a<b$ or $a=b$
or $a>b$ --- note $a<b$ abbreviates $a\leq b$ but $a\neq b$, and $a>b$
abbreviates $b<a$.
\end{definition}

\begin{definition}
Let $(A,\leq)$ be a poset. Let $E\subset A$ be a subset.
An \define{Upper Bound} of $E$ is an element $u\in A$ such that for
all $e\in E$ we have $e\leq u$. Sometimes people write this as
$\exists u\in A\ldotp E\leq u$ in a mild abuse of notation.
\end{definition}

\begin{definition}
Let $(A,\leq)$ be a poset. Let $b\in A$ be an element. We say that $b$
is a \define{Maximal Element} (or just ``maximal'') if for all $a\in A$ 
such that $a\geq b$ we have $a=b$. (Anything not-less-than $b$ must be
equal to $b$.)
\end{definition}

\begin{claim}[Zorn's lemma]
Let $A$ be a \emph{partially ordered set} such that every
\emph{linearly ordered} subset $E\subset A$ has an \emph{upper bound}.
Then $A$ has a \emph{maximal element}.
\end{claim}

\begin{node}
We do not prove this claim, because it is an \emph{axiom} (despite its
name ``Zorn's \emph{lemma}'').
\end{node}

\begin{example}
Let $H$ be a Hilbert space.
Let $\mathcal{A}$ be the set of all orthonormal families in $H$, so we
write $\mathcal{A}=ON(H)$. So $e\in\mathcal{A}$ and $e\subset H$ [so
  $e$ means two things simultaneously].
The ordering is by inclusion $e\leq e'$ iff
$e=\{e_{\alpha}\}_{\alpha\in A}$, $e'=\{e'_{\beta}\}_{\beta\in B}$,
$e\subset e'$.

Let $E\subset\mathcal{A}$ be a linearly ordered subset. Then there
exists an upper bound for $E$, namely $\bigcup E$ since
\begin{equation}
\forall e\in E\ldotp e\leq\bigcup E.
\end{equation}
Therefore by Zorn's lemma, $\mathcal{A}$ has a maximal element. In
particular, this means every Hilbert space must have an orthonormal basis.
\end{example}

\begin{note}
We will use Zorn's Lemma for Hahn--Banach theorem, for Tychanov
theorem, Krane--Miln theorem [for convex sets]. It is used everywhere,
so get used to it.
\end{note}